\subsection{The Dual Problem (DP)}
The primal problem involves minimizing a regularized hinge loss over the parameters defining a separating hyperplane. The dual problem, derived via Lagrangian duality, provides computational and theoretical advantages, especially when expressed in terms of inner products between data points.

The dual formulation offers several key advantages. In the context of support vector machines, 
solving the dual allows the optimization to depend only on the number of data points rather than 
the dimension of the spatial space, making it efficient in high-dimensional settings. 
Additionally, the dual formulation naturally incorporates the kernel trick, enabling non-linear 
decision boundaries without explicitly mapping data to higher dimensions. It also gives insight 
into which data points are support vectors, since only the corresponding dual variables are non-zero. 
Thus, the dual problem is usually more numerically stable and computationally efficient, than the
primal formulation.

\subsubsection{Dual Derivation}
The dual problem is an equivalent formulation of the primal optimization problem derived using Lagrangian multipliers and KKT conditions. It can be expressed as:

\begin{theorem}{The Dual Problem}{dual-problem}
	The dual of the primal problem~\eqref{eq:primal-problem} is given by:
	\[
		\min_{\symbf{\alpha}} \frac{1}{2} \inner*{\symbf{\alpha}, \mathrm{Y} \mathrm{G} \mathrm{Y} \symbf{\alpha}} -  \inner*{\symbf{1}_M, \symbf{\alpha}} \text{ subject to }
		\begin{cases}
			0 \leq \symbf{\alpha} \leq C\symbf{1}_M \\
			\symbf{\alpha}^\top Y = 0
		\end{cases} \tag{DP}
	\]
	Where $G = \bigl(\inner*{\symbf{x}_i, \symbf{x}_j}\bigr)_{i,j=1,\hdots,M} = X X^\top$ is the Gram matrix, and $Y=\operatorname{diag}(y_1,\dots,y_M)$.
\end{theorem}

\begin{proof}{}{}
	\paragraph{Primal problem (P)}
	We begin with the soft-margin SVM primal:
	\[
		\min_{\symbf{w}, b, \symbf{\xi}}
		\tfrac12\|\symbf{w}\|_2^2
		+ C\symbf{1}_M^\top \symbf{\xi}
		\quad\text{subject to}\quad
		\begin{cases}
			c_1(\symbf{w}, b, \symbf{\xi}) &= \symbf{1}_M - Y\left(X\symbf{w} + b\symbf{1}_M\right)-\symbf{\xi}\le\symbf{0}_M, \\
			c_2(\symbf{\xi}) &= -\symbf{\xi}\le\symbf{0}_M.
		\end{cases}
	\]
	\paragraph{Lagrangian}
	We introduce the Lagrange multipliers
	\(\symbf{\alpha}_1,\symbf{\alpha}_2 \ge \symbf{0}_M.
	\)
	The Lagrangian is
	\[
		\mathcal{L}\left(\symbf{w}, b, \symbf{\xi}, \symbf{\alpha}_1,\symbf{\alpha}_2\right) =\tfrac12\|\symbf{w}\|_2^2 +C\symbf{1}_M^\top \symbf{\xi} +\symbf{\alpha}_1^\top c_1(\symbf{w},b,\symbf{\xi}) +\symbf{\alpha}_2^\top c_2(\symbf{\xi}).
	\]

	Because $\mathcal{L}$ is convex in $(\symbf{w},b,\symbf{\xi})$ and concave in $(\symbf{\alpha}_1,\symbf{\alpha}_2)$, the minimax theorem~\cite{vonNeumann1928} implies

	\begin{align*}
		\min_{\symbf{w},b,\symbf{\xi}}\max_{\symbf{\alpha}_1,\symbf{\alpha}_2}\mathcal{L}(\symbf{w},b,\symbf{\xi},\symbf{\alpha}_1,\symbf{\alpha}_2)= \max_{\symbf{\alpha}_1,\symbf{\alpha}_2}\min_{\symbf{w},b,\symbf{\xi}}\mathcal{L}(\symbf{w},b,\symbf{\xi},\symbf{\alpha}_1,\symbf{\alpha}_2) = \max_{\symbf{\alpha}_1,\symbf{\alpha}_2} g(\symbf{\alpha}_1,\symbf{\alpha}_2)
	\end{align*}

	\paragraph{KKT conditions}
	To find the minimum of the Lagrangian, we need to satisfy the KKT conditions. The KKT conditions are a set of necessary and sufficient conditions for optimality in constrained optimization problems. They consist of three components:
	\begin{itemize}
		\item \textbf{Stationarity} The gradients of the Lagrangian with respect to the primal variables $(\symbf{w}, b, \symbf{\xi})$ must vanish:
		\begin{align*}
			\nabla_{\symbf{w}} \Lcal = \symbf{w} -  (X^\top Y) \symbf{\alpha}_1 = \symbf{0}, \quad \nabla_{b} \Lcal = - \symbf{\alpha}_1^\top Y = \symbf{0}^\top, \quad \nabla_{\symbf{\xi}} \Lcal  = C \symbf{1}_M^\top - \symbf{\alpha}_1^\top - \symbf{\alpha}_2^\top = \symbf{0}_M^\top
		\end{align*}
		\item \textbf{Complementary slackness} The Lagrange multipliers $\symbf{\alpha}_1$ and $\symbf{\alpha}_2$ must satisfy the complementary slackness conditions:
		\begin{align*}
			\symbf{\alpha}_1^\top c_1(\symbf{w},b,\symbf{\xi}) = \symbf{0}_M^\top, \quad \symbf{\alpha}_2^\top c_2(\symbf{\xi}) = \symbf{0}_M^\top
		\end{align*}
		\item \textbf{Feasibility} The primal variables $(\symbf{w}, b, \symbf{\xi})$ must satisfy the primal constraints:
		\begin{align*}
			c_1(\symbf{w},b,\symbf{\xi}) = \symbf{1}_M - Y\left(X\symbf{w} + b\symbf{1}_M\right)-\symbf{\xi}\le\symbf{0}_M, \\
			c_2(\symbf{\xi}) = -\symbf{\xi}\le\symbf{0}_M.
		\end{align*}
	\end{itemize}
	Thus, we have the following conditions:
	\begin{align*}
		\symbf{w} &= X^\top Y \symbf{\alpha}_1 , \quad \symbf{\alpha}_1^\top Y  = \symbf{0}_M^\top, \quad \symbf{0}_M \leq \symbf{\alpha}_1 + \underset{=\symbf{0}}{\symbf{\alpha}_2} \leq C\symbf{1}_M \; \text{where} \\
		\norm*{\symbf{w}}_2^2 &= \norm{X^\top Y \symbf{\alpha}_1}_2^2 = \left(X^\top Y \symbf{\alpha}_1\right)^\top \left(X^\top Y \symbf{\alpha}_1\right) = \symbf{\alpha}_1^\top Y [X X^\top] Y \symbf{\alpha}_1 = \inner*{\symbf{\alpha}_1, \mathrm{Y} \mathrm{G} \mathrm{Y} \symbf{\alpha}_1}
	\end{align*}
	Substituting this into the Lagrangian gives:
	\begin{align*}
		g(\symbf{\alpha}_1) = \tfrac12\symbf{\alpha}_1^\top YGY \symbf{\alpha}_1 -\symbf{\alpha}_1^\top YGY\symbf{\alpha}_1 +\symbf{\alpha}_1^\top \symbf{1}_M = -\tfrac12\inner*{\symbf{\alpha}_1, \mathrm{Y} \mathrm{G} \mathrm{Y} \symbf{\alpha}_1} + \inner*{\symbf{1}_M, \symbf{\alpha}_1}
	\end{align*}
	\paragraph{Dual problem (DP)}
	Flipping signs for a minimization form, gives us the dual problem is
	\[
		\min_{\symbf{\alpha}} \tfrac12\symbf{\alpha}^\top YGY \symbf{\alpha}
		-\symbf{1}_M^\top \symbf{\alpha}
		\quad\text{subject to}\quad
		0\le\symbf{\alpha}\le C\symbf{1}_M,
		\symbf{\alpha}^\top Y =0.
	\]
	This completes the derivation of the dual problem~\eqref{eq:dual-problem}.\qed
\end{proof}

Once the dual problem \eqref{eq:dual-problem} has been solved and an optimal dual variable 
$\symbf{\alpha}^\star$ has been obtained, we can reconstruct the corresponding solution 
$(\symbf{w}^\star, b^\star, \symbf{\xi}^\star)$ to the primal problem \eqref{eq:primal-problem} 
using the KKT optimality conditions.
\begin{proposition}{Recovering Primal Solution from Dual Solution}{D2P-solution}
	Given an optimal solution $\symbf{\alpha}^\star$ to the dual problem, the corresponding optimal solution $(\symbf{w}^\star, b^\star, \symbf{\xi}^\star)$ to the primal problem can be recovered through:
	\begin{align*}
		\symbf{w}^\star &= X^\top Y \symbf{\alpha}^\star \\
		b^\star &= y_i - \symbf{x}_i^\top\symbf{w}^\star \quad \text{for any index } i \text{ such that } 0 < \alpha_i^\star < C \\
		\symbf{\xi}^\star &= \max\left\{\symbf{0}_M,\; \symbf{1}_M - Y(X\symbf{w}^\star + b^\star\symbf{1}_M)\right\}
	\end{align*}
\end{proposition}

\subsubsection{Existence and Uniqueness of the Dual Solution}
The dual problem is a convex quadratic program with linear constraints. 
The existence and uniqueness of the solution depend on the properties of the Gram matrix $G$ and the constraints imposed by the labels $Y$.

\begin{theorem}{Existence \& Uniqueness of Dual Solution}{existence-uniqueness-dual}
	The dual problem \eqref{eq:dual-problem} always admits at least one global minimizer. Moreover, if $YGY$ is strictly positive-definite on the subspace $\{\symbf{\alpha}\mid \inner{\symbf{y}, \symbf{\alpha}}=0\}$, then the solution $\symbf{\alpha}^\star$ is unique. In degenerate cases, multiple optimal solutions may exist.
\end{theorem}

\begin{proof}{}{}
	\paragraph{Existence of a solution:}
	The dual problem \eqref{eq:dual-problem} is a convex quadratic program in $\symbf{\alpha}$, subject to a linear equality constraint $\symbf{\alpha}^\top Y = \symbf{0}^\top$ and box constraints $0 \leq \symbf{\alpha} \leq C\symbf{1}_M$. We define the feasible region as:
	\[
		\Omega = \left\{ \symbf{\alpha} \in \mathbb{R}^M \;\middle|\; 0 \leq \symbf{\alpha} \leq C\symbf{1}_M, \;\symbf{\alpha}^\top Y = \symbf{0}^\top \right\},
	\]
	which is the intersection of the closed box $[0, C]^M$ with a hyperplane. This intersection is nonempty: because the dataset includes both positive and negative labels (i.e., $\exists\,p,n$ such that $y_p = 1$, $y_n = -1$), one can construct a feasible point by choosing $\alpha_p = \alpha_n$ and setting all other components to zero. Hence, $\Omega$ is closed, bounded, and nonempty.

	The dual objective function,
	\[
		f(\symbf{\alpha}) = \frac{1}{2} \inner*{\symbf{\alpha}, \mathrm{Y} \mathrm{G} \mathrm{Y} \symbf{\alpha}} - \inner*{\symbf{1}_M, \symbf{\alpha}},
	\]
	is continuous and convex, being a quadratic form with symmetric matrix $YGY$ and linear term $-\symbf{1}_M^\top \symbf{\alpha}$. Since $f$ is continuous and convex over a compact feasible set $\Omega$, the extreme value theorem guarantees the existence of a global minimizer $\symbf{\alpha}^\star$.

	\paragraph{Uniqueness of a solution:}
	The Hessian of $f$ is $\nabla^2 f(\symbf{\alpha}) = YGY$, which is symmetric and positive semi-definite because $G$ is a Gram matrix and $Y$ is a diagonal matrix with $\pm 1$ entries. 
	If $YGY$ is strictly positive-definite on the subspace $\{\symbf{\alpha} \in \mathbb{R}^M : \symbf{\alpha}^\top Y = \symbf{0}^\top\}$, then $f$ is strictly convex on the feasible region $\Omega$, and the minimizer $\symbf{\alpha}^\star$ is unique. 
	In contrast, if $YGY$ is only positive semi-definite (e.g., if $G$ is rank-deficient), then flat directions may exist in the objective, and multiple global minimizers may be possible.

	\paragraph{Degenerate cases:}
	Degenerate cases occur when $YGY$ is positive semi-definite but not positive-definite on the constraint subspace. This happens in several scenarios:
	\begin{enumerate}
	\item When the data vectors $\{\symbf{x}_i\}$ are linearly dependent, making the Gram matrix $G$ rank-deficient. For example, if two feature vectors are identical but have the same label, or if one feature vector is a linear combination of others.
	\item When the number of data points $M$ exceeds the dimension $N$ of the feature space, necessarily creating linear dependencies.
	\end{enumerate}
	In these cases, there exist non-zero directions $\symbf{v}$ satisfying $\symbf{v}^\top YGY \symbf{v} = 0$ and $\symbf{v}^\top Y = 0$. Moving along such directions from one optimal solution leads to another optimal solution, creating a flat region (typically a face or edge) of the feasible region where all points are global minimizers. Although $\symbf{w}^\star$ remains unique due to the form $\symbf{w}^\star = X^\top Y \symbf{\alpha}^\star$, the values of $\symbf{\alpha}^\star$ corresponding to this unique $\symbf{w}^\star$ are not unique.
\end{proof}
